\documentclass[usenames,dvipsnames, 18pt, compress, aspectratio=169]{beamer}

% can be compiled by xelatex -shell-escape presentation.tex
% lualatex -shell-escape presentation.tex

\usetheme[]{metropolis}

\usepackage[utf8]{inputenc}
\usepackage[russian, english]{babel}
\usepackage{booktabs}
\usepackage[scale=2]{ccicons}
\usepackage{listings}
\usepackage{marvosym}
\usepackage{color}
\usepackage{xcolor}
\usepackage[document]{ragged2e}
\usepackage[export]{adjustbox}
\usepackage{fontawesome5}
\usepackage{enumitem}
\usepackage{minted}
\usemintedstyle{tango}
\usepackage[normalem]{ulem}
\usepackage{tikz}
\usetikzlibrary{patterns}
\usetikzlibrary{mindmap}
\usetikzlibrary{shapes.misc, fit}
\usetikzlibrary{spy, decorations, decorations.pathmorphing, decorations.pathreplacing, backgrounds, decorations.text}
\usepackage{graphicx}
\usepackage{eso-pic}
\usepackage{verbatim}
\usepackage{smartdiagram}
\usesmartdiagramlibrary{additions}
\usetikzlibrary{trees}
\usepackage{datetime}
\usepackage{hyperref}
\usepackage{forloop}
\usepackage{csquotes}
\usetikzlibrary{tikzmark}
\usetikzlibrary{arrows.meta}

\usepackage{tcolorbox}
\usepackage{tabularx}
\usepackage{array}
\usepackage{colortbl}
\tcbuselibrary{skins}

\usetikzlibrary{shapes,arrows,positioning}
\graphicspath{{images/}}
%\newfontfamily{\FA}{FontAwesome}

\def\twitter{{\faTwitter}}
\def\github{{\faGithub}}
\def\email{{\faEnvelope}}

\renewcommand{\ttdefault}{pcr}
\setmonofont{FiraCode-VF}
%\newfontfamily{\ttfamily}{FiraCode-VF}

\usefonttheme{professionalfonts} % using non standard fonts for beamer
\usefonttheme{serif} % default family is serif
\usepackage{fontspec}
\setmainfont{Liberation Sans}
\newfontfamily\ExtraLight{Liberation Sans}
\newfontfamily\Light{Liberation Sans}
\newfontfamily\Book{Liberation Sans}
\newfontfamily\Medium{Liberation Sans}

\makeatletter
\newcommand\HUGE{\@setfontsize\Huge{32}{41}}
\makeatother

\newcommand\AtPagemyUpperLeft[1]{\AtPageLowerLeft{%
\put(\LenToUnit{0.85\paperwidth},\LenToUnit{0.05\paperheight}){#1}}}

\newcommand\AtPagemyUpperTop[1]{\AtPageLowerLeft{%
\put(\LenToUnit{0.42\paperwidth},\LenToUnit{0.90\paperheight}){#1}}}

\renewcommand{\ULthickness}{2.0pt}

\definecolor{links}{HTML}{0099FF}
\hypersetup{colorlinks, linkcolor=, urlcolor=links}
\definecolor{greenGood}{HTML}{99FF99}
\definecolor{redBad}{HTML}{FF9980}
\definecolor{title}{HTML}{ee0000}

\setbeamerfont{section title}{family=\Book, size=\Huge, shape=\normalfont}
\setbeamerfont{frametitle}{family=\Book, size=\large, shape=\normalfont}
\setbeamerfont{title}{family=\Book, size=\Large, shape=\normalfont}
\setbeamerfont{subtitle}{size=\small}
\setbeamerfont{author}{family=\ExtraLight, size=\footnotesize}

\definecolor{cec1d24}{RGB}{236,29,36}
\definecolor{cffffff}{RGB}{255,255,255}

\setbeamertemplate{navigation symbols}{}
\beamertemplatenavigationsymbolsempty
\pagenumbering{gobble}

\setbeamertemplate{title page}
{

  \vspace*{2.1cm}
  \hspace{7.0cm}
  \begin{minipage}[b][\paperheight]{0.5\textwidth}
  \begin{center}

    \ifx\inserttitle\@empty\else
    {{% \inserttitle is nonempty
      \raggedright%
      %\linespread{1.0}%
      \usebeamerfont{title}%
      \usebeamercolor[fg]{title}%
      %\vspace*{1.3em}
      \if@noSmallCapitals%
        \inserttitle%
      \else%
        \scshape{\color{title} \textbf{\begin{flushleft}\inserttitle\end{flushleft}}}%
      \fi%
      \vspace*{0.3em}
    }}
    \fi

    \vspace*{0.5em}%

    \ifx\insertsubtitle\@empty\else
    {{% \insertsubtitle is nonempty
      \usebeamerfont{subtitle}%
      \usebeamercolor[fg]{subtitle}%
      {\color{black} \insertsubtitle}%
      \vspace*{3.0em}%
    }}
    \fi

    \vspace*{1.0em}%

    \usebeamerfont{author}%
    \usebeamercolor[fg]{author}%
    {\begin{flushleft}\color{black} \insertauthor\end{flushleft}}%

    %\vspace*{1.5em}
    \fontsize{8pt}{10}\selectfont
    {\begin{flushleft}\color{black} 24-09-2022\end{flushleft}}%

    \vfill
    \vspace*{2em}
  \end{center}
  \end{minipage}
}

\setbeamertemplate{section page}
{
  \vspace{2em}
  \centering
  \begin{minipage}{22em}
    \usebeamercolor[fg]{section title}
    \usebeamerfont{section title}
    {\color{black} \insertsectionhead\\[-1ex]}
  \end{minipage}
  \par
}

\setbeamertemplate{footline}
{
\begin{beamercolorbox}[wd=\textwidth,ht=3ex,dp=3ex,leftskip=0.3cm,rightskip=0.3cm]{structure}
  \usebeamerfont{page number in head/foot}
  \insertframenumber
\end{beamercolorbox}
}


\title{Party tricks for PostgreSQL \\ perf, ftrace and bpftrace}
\subtitle{}
\date{\today}
\author{DMITRII DOLGOV}
\institute{}

\begin{document}
{
  \usebackgroundtemplate{\includegraphics[width=\paperwidth]{title.png}}%
  \fontsize{17pt}{18}\selectfont
  \maketitle
}

\AddToShipoutPictureBG{
  \AtPagemyUpperLeft{{\includegraphics[width=2.0cm,keepaspectratio]{logo.png}}}
}

%\setbeamertemplate{background canvas}{}

\fontsize{17pt}{18}\selectfont

{
\setbeamercolor{background canvas}{bg=black}
\begin{frame}
    \frametitle{}
    \begin{center}
        \includegraphics[width=1.0\textwidth,center]{nethack.png}
    \end{center}
\end{frame}
\note{Originally about resources, but now not only}
}

\setbeamertemplate{background canvas}{}

\begin{frame}[fragile]{}
    \frametitle{}
    \begin{center}
        \begin{minted}[fontsize=\large]{c}
    /* XXX guessed-at value */
    #define PGPROC_MAX_CACHED_SUBXIDS 64


    /* has PGPROC->subxids overflowed */
    bool        overflowed;
        \end{minted}
    \end{center}
\end{frame}
\note{
    pg_stat_toast, original FPI stats (have landed in pgss via wal_fpi), etc.
}

\begin{frame}[fragile]{}
    \frametitle{}
    \begin{center}
        \begin{minted}[fontsize=\large]{shell-session}
$ perf record -e probe_postgres:GetNewTransactionId
$ perf script
    postgres 29959 probe_postgres
    GetNewTransactionId: nxids=0

        \end{minted}
    \end{center}
\end{frame}
\note{edited for slides}

\section{Instrumentation basics}

\begin{frame}[fragile]{}
    \frametitle{}
    \begin{center}
        \begin{itemize}[label={}]
            \item {Static tracing}
            \item {Tracepoints}
            \item {USDT}
        \end{itemize}
    \end{center}
\end{frame}

\begin{frame}[fragile]{}
    \frametitle{}
    \begin{center}
        \begin{itemize}[label={}]
            \item {Dynamic tracing}
            \item {kprobe/kretprobe}
            \item {uprobe/uretprobe}
        \end{itemize}
    \end{center}
\end{frame}

\begin{frame}[fragile]{}
    \frametitle{}
    \begin{center}
        \begin{minted}[fontsize=\Large, escapeinside=||]{c}
    <handle_event>:   |\colorbox{red!20}{push %rbp}|
    <handle_event+1>: mov  %rsp,%rbp
    <handle_event+4>: sub  $0x60,%rsp

    <handle_event>:   |\colorbox{red!20}{int3}|
    <handle_event+1>: mov  %rsp,%rbp
    <handle_event+4>: sub  $0x60,%rsp
        \end{minted}

    \end{center}
\end{frame}
\note{
    cat /proc/30615/maps
    00400000-0045a000 r-xp 00000000 fd:03 338992387 /path/to/binary

    gdb -p <process>

    > x/128i 0x00400000

    Breakpoint/context switch overhead.
}

\section{Static tracing}

\begin{frame}[fragile]{}
    \frametitle{}
    \begin{center}
        \begin{minted}[fontsize=\large]{shell-session}
--enable-dtrace

$ bpftrace -l 'usdt:bin/postgres:*'
usdt:bin/postgres:postgresql:buffer__checkpoint__done
usdt:bin/postgres:postgresql:buffer__checkpoint__start
[...]
        \end{minted}
    \end{center}
\end{frame}
\note{edited for slides}

\begin{frame}[fragile]{}
    \frametitle{}
    \begin{center}
        \begin{minted}[fontsize=\large]{shell-session}
$ bpftrace -p 31715 -f json -e
    'usdt:bin/postgres:postgresql:checkpoint__start
    {printf("Checkpoint, Flags %d\n", arg0);}'

{"type": "attached_probes", "data": {"probes": 1}}
{"type": "printf", "data": "Checkpoint, Flags 108\n"}
{"type": "printf", "data": "Checkpoint, Flags 108\n"}
        \end{minted}
    \end{center}
\end{frame}
\note{edited for slides}

\section{Dynamic tracing}

\begin{frame}[fragile]{}
    \frametitle{}
    \begin{center}
        \begin{minted}[fontsize=\large]{shell-session}
$ perf probe -x bin/postgres \
    -L GetNewTransactionId

[...]
187 int nxids = MyProc->subxidStatus.count;
[...]
        \end{minted}
    \end{center}
\end{frame}

\begin{frame}[fragile]{}
    \frametitle{}
    \begin{center}
        \begin{minted}[fontsize=\large]{shell-session}
$ perf probe -x bin/postgres \
    -V GetNewTransactionId:187

Available variables at GetNewTransactionId:187
	@<GetNewTransactionId+1040>
		FullTransactionId       full_xid
		TransactionId   xid
		XidCacheStatus* substat
		_Bool   isSubXact
		char*   __func__
		int     nxids
        \end{minted}
    \end{center}
\end{frame}

\begin{frame}[fragile]{}
    \frametitle{}
    \begin{center}
        \begin{minted}[fontsize=\large]{shell-session}
$ perf probe -x bin/postgres \
    GetNewTransactionId:187 nxids

$ perf record -e \
	probe_postgres:GetNewTransactionId_L187_2

$ perf script
	postgres probe_postgres
	GetNewTransactionId_L187_2: nxids=0

        \end{minted}
    \end{center}
\end{frame}
\note{edited for slides}

\section{Limitations}

\begin{frame}[fragile]{}
    \frametitle{}
    \begin{center}
        \begin{minted}[fontsize=\large, escapeinside=::]{c}
*dovacuum =
	force_vacuum ||
	(vactuples > vacthresh) ||
	(vac_ins_base_thresh >= 0 &&
		:\colorbox{red!20}{instuples > vacinsthresh)};:
        \end{minted}
    \end{center}
\end{frame}
\note{
	Filter doesn't work with comparisons, requires post-processing,
	complicated. For such cases it's good to have a stable usdt.
}

\begin{frame}[fragile]{}
    \frametitle{}
    \begin{center}
        \begin{itemize}[label={}]
            \item<+-> {Context switch overhead}
            \item<+-> {Debugging symbols}
            \item<+-> {Incorporate into infrastructure?}
            \item<+-> {Privileges?}
        \end{itemize}
    \end{center}
\end{frame}

\begin{frame}[fragile]{}
    \frametitle{}
    \begin{center}
        \begin{minted}[fontsize=\large]{c}
BPF_OBJ_PIN
    Pin an eBPF program or map referred
    by the specified *bpf_fd* to the
    provided *pathname* on the filesystem.
        \end{minted}
    \end{center}
\end{frame}
\note{
    BPFD
}

\begin{frame}[fragile]{}
    \frametitle{}
    \begin{center}
        \begin{minted}[fontsize=\large]{shell-session}
$ cat events/synthetic/query_latency/hist
# event histogram

hitcount: 1
{ lat_us: 223,
  query_string: select a from test1;}

hitcount: 1
{ lat_us: 314,
  query_string: select b from test2;}
        \end{minted}
    \end{center}
\end{frame}

\begin{frame}[fragile]{}
    \frametitle{}
    \begin{center}
        \begin{minted}[fontsize=\large]{shell-session}
$ perf probe -x bin/postgres \
    exec_simple_query query_string:string

$ perf probe -x bin/postgres \
    exec_simple_query%return
        \end{minted}
    \end{center}
\end{frame}

\begin{frame}
    \frametitle{}
    \begin{center}
        \includegraphics[width=0.9\textwidth,center]{more.png}
    \end{center}
\end{frame}

\begin{frame}[fragile]{}
    \frametitle{}
    \begin{center}
    \begin{overprint}[14.5cm]
        \onslide<1>
		\vspace{-0.73cm}
        \begin{minted}[fontsize=\large, bgcolor=white]{shell-session}
synthetic_events
query_latency u64 lat_us; char[] query_string
        \end{minted}

		\vspace{-0.73cm}
        \begin{minted}[fontsize=\large, bgcolor=white]{shell-session}
events/probe_postgres/exec_simple_query/trigger
hist:keys=common_pid:t0=common_timestamp.usecs
        \end{minted}

		\vspace{-0.73cm}
        \begin{minted}[fontsize=\large, bgcolor=white]{shell-session}
events/probe_postgres/exec_simple_query__return/trigger
hist:keys=common_pid:lat_us=common_timestamp.usecs-$t0:
   onmatch(probe_postgres.exec_simple_query).
   trace(query_latency,$lat_us,query_string_string)
        \end{minted}

		\vspace{-0.73cm}
        \begin{minted}[fontsize=\large, bgcolor=white]{shell-session}
events/synthetic/query_latency/trigger
hist:keys=lat_us,query_string:sort=lat_us
        \end{minted}

        \onslide<2>
        \begin{minted}[fontsize=\large, bgcolor=red!20]{shell-session}
synthetic_events
query_latency u64 lat_us; char[] query_string
        \end{minted}

		\vspace{-0.73cm}
        \begin{minted}[fontsize=\large, bgcolor=white]{shell-session}
events/probe_postgres/exec_simple_query/trigger
hist:keys=common_pid:t0=common_timestamp.usecs
        \end{minted}

		\vspace{-0.73cm}
        \begin{minted}[fontsize=\large, bgcolor=white]{shell-session}
events/probe_postgres/exec_simple_query__return/trigger
hist:keys=common_pid:lat_us=common_timestamp.usecs-$t0:
   onmatch(probe_postgres.exec_simple_query).
   trace(query_latency,$lat_us,query_string_string)
        \end{minted}

		\vspace{-0.73cm}
        \begin{minted}[fontsize=\large, bgcolor=white]{shell-session}
events/synthetic/query_latency/trigger
hist:keys=lat_us,query_string:sort=lat_us
        \end{minted}

        \onslide<3>
        \begin{minted}[fontsize=\large, bgcolor=white]{shell-session}
synthetic_events
query_latency u64 lat_us; char[] query_string
        \end{minted}

		\vspace{-0.73cm}
        \begin{minted}[fontsize=\large, bgcolor=red!20]{shell-session}
events/probe_postgres/exec_simple_query/trigger
hist:keys=common_pid:t0=common_timestamp.usecs
        \end{minted}

		\vspace{-0.73cm}
        \begin{minted}[fontsize=\large, bgcolor=white]{shell-session}
events/probe_postgres/exec_simple_query__return/trigger
hist:keys=common_pid:lat_us=common_timestamp.usecs-$t0:
   onmatch(probe_postgres.exec_simple_query).
   trace(query_latency,$lat_us,query_string_string)
        \end{minted}

		\vspace{-0.73cm}
        \begin{minted}[fontsize=\large, bgcolor=white]{shell-session}
events/synthetic/query_latency/trigger
hist:keys=lat_us,query_string:sort=lat_us
        \end{minted}

        \onslide<4>
        \begin{minted}[fontsize=\large, bgcolor=white]{shell-session}
synthetic_events
query_latency u64 lat_us; char[] query_string
        \end{minted}

		\vspace{-0.73cm}
        \begin{minted}[fontsize=\large, bgcolor=white]{shell-session}
events/probe_postgres/exec_simple_query/trigger
hist:keys=common_pid:t0=common_timestamp.usecs
        \end{minted}

		\vspace{-0.73cm}
        \begin{minted}[fontsize=\large, bgcolor=red!20]{shell-session}
events/probe_postgres/exec_simple_query__return/trigger
hist:keys=common_pid:lat_us=common_timestamp.usecs-$t0:
   onmatch(probe_postgres.exec_simple_query).
   trace(query_latency,$lat_us,query_string_string)
        \end{minted}

		\vspace{-0.73cm}
        \begin{minted}[fontsize=\large, bgcolor=white]{shell-session}
events/synthetic/query_latency/trigger
hist:keys=lat_us,query_string:sort=lat_us
        \end{minted}

        \onslide<5>
        \begin{minted}[fontsize=\large, bgcolor=white]{shell-session}
synthetic_events
query_latency u64 lat_us; char[] query_string
        \end{minted}

		\vspace{-0.73cm}
        \begin{minted}[fontsize=\large, bgcolor=white]{shell-session}
events/probe_postgres/exec_simple_query/trigger
hist:keys=common_pid:t0=common_timestamp.usecs
        \end{minted}

		\vspace{-0.73cm}
        \begin{minted}[fontsize=\large, bgcolor=white]{shell-session}
events/probe_postgres/exec_simple_query__return/trigger
hist:keys=common_pid:lat_us=common_timestamp.usecs-$t0:
   onmatch(probe_postgres.exec_simple_query).
   trace(query_latency,$lat_us,query_string_string)
        \end{minted}

		\vspace{-0.73cm}
        \begin{minted}[fontsize=\large, bgcolor=red!20]{shell-session}
events/synthetic/query_latency/trigger
hist:keys=lat_us,query_string:sort=lat_us
        \end{minted}

    \end{overprint}

    \end{center}
\end{frame}

\section{Development}

\begin{frame}[fragile]{}
    \frametitle{}
    \begin{center}
        \begin{minted}[fontsize=\large]{shell-session}
$ perf record -a -g --event=mem:0x7ff3bc1f9000/8:rw
        \end{minted}
    \end{center}
\end{frame}
\note{share memory from postmaster maps}

\begin{frame}[fragile]{}
    \frametitle{}
    \begin{center}
        \begin{minted}[fontsize=\large]{text}
$ bpftrace -e 'kprobe:__x64_sys_pwrite64
    /comm == "postgres"/ { override(-12) }' --unsafe
Attaching 1 probe...
        \end{minted}

        \begin{minted}[fontsize=\large]{text}
=# create table test();
PANIC:  XX000: could not write to
log file 00000001000000000000002B at offset 11444224,
length 106496: Cannot allocate memory
LOCATION:  XLogWrite, xlog.c:2233
        \end{minted}
    \end{center}
\end{frame}
\note{ENOMEM, only for those functions that are markes safe to fail.}

\section{Tips \& tricks}

\begin{frame}[fragile]{}
    \frametitle{}
    \begin{center}
        \begin{minted}[fontsize=\large]{shell-session}
$ perf report --header

# ========
# captured on    : Wed Oct 12 17:07:32 2022
# header version : 1
# data offset    : 504
# data size      : 1392968
# feat offset    : 1393472
# arch : x86_64
# cpudesc : Intel(R) Core(TM) i7-10610U CPU @ 1.80GHz
# ...
        \end{minted}
    \end{center}
\end{frame}
\note{ENOMEM, only for those functions that are markes safe to fail.}

\begin{frame}[fragile]{}
    \frametitle{}
    \begin{center}
        \begin{minted}[fontsize=\large]{shell-session}
$ perf archive

  Now please run:
  tar xvf perf.data.tar.bz2 -C ~/.debug
        \end{minted}
    \end{center}
\end{frame}
\note{Archives all available symbols}

\begin{frame}[fragile]{}
    \frametitle{}
    \begin{center}
        \begin{minted}[fontsize=\large]{shell-session}
$ perf record --overwrite --switch-output=signal
$ kill -SIGUSR2 <perf-pid>
        \end{minted}
    \end{center}
\end{frame}
\note{Daemonized perf, optionally specify ring buffer size with -m 1M to
    control memory better}

\begin{frame}[fragile]{}
    \frametitle{}
    \begin{center}
        \begin{minted}[fontsize=\large]{shell-session}
$ perf record -p <pid> -e intel_pt//u --filter\
 'filter tts_buffer_heap_materialize @ bin/postgres'
        \end{minted}
    \end{center}
\end{frame}
\note{Duplicated functions?}

\begin{frame}
    \frametitle{}
    \begin{center}
		\begin{tikzpicture}[
		  decoration={
			reverse path,
			text effects along path,
			text={perf record -e probe\_postgres:GetNewTransactionId perf script bpftrace -l 'usdt:bin/postgres:*' bpftrace -p 31715 -f json -e 'usdt:bin/postgres:postgresql:checkpoint\_\_start printf("Checkpoint, Flags d", arg0);' perf probe -x bin/postgres -L GetNewTransactionId perf probe -x bin/postgres -V GetNewTransactionId:187 perf probe -x bin/postgres GetNewTransactionId:187 nxids perf record -e probe\_postgres:GetNewTransactionId\_L187\_2 perf script cat events/synthetic/query\_latency/hist perf probe -x bin/postgres exec\_simple\_query query\_string:string perf probe -x bin/postgres exec\_simple\_query return echo 'query\_latency u64 lat\_us; char[] query\_string' > synthetic\_events echo 'hist:keys=common\_pid:ts0=common\_timestamp.usecs' > events/probe\_postgres/exec\_simple\_query/trigger echo 'hist:keys=common\_pid:lat\_us=common\_timestamp.usecs-ts0:onmatch(probe\_postgres.exec\_simple\_query).trace(query\_latency,lat\_us,query\_string\_string)' > events/probe\_postgres/exec\_simple\_query\_\_return/trigger echo 'hist:keys=lat\_us,query\_string:sort=lat\_us' > events/synthetic/query\_latency/trigger perf record -a -g --event=mem:0x7ff3bc1f9000/8:rw bpftrace -e 'kprobe:\_\_x64\_sys\_pwrite64 comm == "postgres"  override(-12) ' --unsafe perf report --header perf archive perf record --overwrite --switch-output=signal kill -SIGUSR2 <perf-pid> perf record -p <pid> -e intel\_pt//u --filter 'filter tts\_buffer\_heap\_materialize @ bin/postgres'},
			text effects/.cd,
			  text along path,
			  character count=\i, character total=\n,
			  characters={scale=1-\i/\n}
			}
		]
		\draw [decorate] (0,0)
			\foreach \i [evaluate={\r=(\i/2000)^2;}] in {0,5,...,4880}{ -- (\i:\r)};
		\end{tikzpicture}
    \end{center}
\end{frame}

\begin{frame}
    \frametitle{}
    \begin{center}
        \includegraphics[width=0.7\textwidth,center]{elephant.jpg}
    \end{center}
\end{frame}

\fontsize{18pt}{18}\selectfont
\begin{frame}
  \vspace*{2.5cm}
  \begin{minipage}[b][\paperheight]{\textwidth}
  \begin{center}

      %\raggedright%
      \linespread{1.0}%
      \usebeamerfont{title}%
      \usebeamercolor[fg]{title}%
      \if@noSmallCapitals%
        Questions?
      \else%
        \scshape{\color{black} Questions?}%
      \fi%
      \vspace*{0.3em}

      \usebeamerfont{subtitle}%
      \fontsize{13pt}{14}\selectfont
      \usebeamercolor[fg]{subtitle}%
        \begin{itemize}[label={}]
            \item {\color{black} \twitter\ @erthalion}
            \item {\color{black} \email\ ddolgov at redhat dot com}
        \end{itemize}
      \vspace*{2.5em}%

    \vfill
    \vspace*{2em}
  \end{center}
  \end{minipage}

\end{frame}

\end{document}
